\documentclass[11pt]{article}

\usepackage{amsmath,amssymb,amstext,amsthm}
\usepackage{geometry}
\usepackage{fancyhdr}
\usepackage[pdftex]{graphicx}
\usepackage{xcolor}

\geometry{
  verbose,
  dvips,
  width=430pt, marginparsep=0pt, marginparwidth=0pt,
  top=90pt, headheight=12pt, headsep=20pt, footskip=30pt, bottom=80pt}

\setlength{\parskip}{\medskipamount}
\setlength{\parindent}{0pt}

\linespread{1.05}

\newtheorem{theorem}{Theorem}
\newtheorem{lemma}[theorem]{Lemma}
\newtheorem{proposition}[theorem]{Proposition}
\newtheorem{corollary}[theorem]{Corollary}
\newtheorem{conjecture}[theorem]{Conjecture}
\newtheorem{obs}{Observation}
\newtheorem{clm}{Claim}
\newtheorem{ques}{Question}
\newtheorem{problem}{Problem}

\newtheorem{ex}{Example}


\newcommand{\spc}{\hspace{0.08in}}
\newcommand{\vs}{\vspace*{.1in}}
\newcommand{\E}{\mathbb{E}}
\newcommand{\Prob}{\mathbb{P}}
\newcommand{\edit}[1]{\emph{\textcolor{blue}{#1}}}


\renewenvironment{proof}{{\noindent \textbf{\textit{Proof.}}}}
{\hfill $\Box$ \vspace*{0.1in}}
\newenvironment{proofc}{{\noindent \textit{Proof of Claim.}}}
{\hfill $\Box$}


\begin{document}
\begin{huge}
    \begin{center}
        Test 2 Review
    \end{center}
\end{huge}


\vs\vs

\section{Limits}

\begin{problem}
Evaluate the following limit. If it does not exists write DNE. \[\lim_{(x,y,z) \to (-1,1,0)} e^{xyz}+x^2\cos{z}-2x\tan^{-1}{y}\]
\end{problem}

\begin{problem}
 Use the Squeeze Theorem to Evaluate the following limit. \[\lim_{(x,y) \to (0,0)} x^4\sin\left(\frac{1}{x^2+y^2}\right)\]   
\end{problem}

\begin{problem}
Evaluate the following limit. If it does not exists write DNE. \[\lim_{(x,y) \to (2,2)} \frac{x-y}{\sqrt{x}-\sqrt{y}}\]
\end{problem}

\begin{problem}
    Evaluate the following limit. If it does not exist write DNE. \[\lim_{(x,y)\to (2,0)} \frac{2xy-2y}{x^2+y^2-4x+4}\]
\end{problem}

\begin{problem}
Evaluate the following limit. If it does not exists write DNE. \[\lim_{(x,y) \to (0,0)} \frac{x^6-3y^2x^3}{y^4+2x^6}\]
\end{problem}


\section{Partial Derivatives}

\begin{problem}
Find both first order partial derivatives of the following surface.
\begin{enumerate}
    \item $f(x,y) = 4^{\cos(xy)}+2xy^5$
    \item $f(x,y) = x^2\sqrt[6]{2xy}$
    \item $z = \frac{6y}{\ln{(x+y)}}$
    \item $f(x,y) = \cos^2{(ye^{x})}$
\end{enumerate}
\end{problem}

\begin{problem}
    Find all second order partial 
\begin{enumerate}
    \item $f(x,y) = \frac{x^3}{y^5}+x^2-y^7$
    \item $f(x,y) = x^{\cos{y}}+e^x$
\end{enumerate}
\end{problem}

\begin{problem}
Consider the surface $f(x,y) = 5x^3y-2y^2-3x+1$. Find the direction of steepest ascent at $(2,2,67)$ and find the slope of steepest ascent.
\end{problem}




\section{The Chain Rule}

\begin{problem}
Find $\frac{\partial z}{\partial t}$ if $z = \sqrt{x}-\ln{xy}$, $x = st+t^2$, and $y = \frac{5s}{t}$.
\end{problem}

\begin{problem}
    Find $\frac{df}{d\theta}$ evaluated at $\theta = \pi$ where $f(x,y) = e^{xy^2}$, $x = \cos{\theta}-\sin^2{\theta}$ and $y = \tan{\left(\frac{\theta}{4}\right)}$.
\end{problem}

\begin{problem}
    Find an equation which represents $\frac{\partial z}{\partial m}$ where $z = f_1(x,y,z)$, $x = f_2(s)$, $y = f_3(t)$, $z = f_4(s,t)$, $s = f_5(n)$, $t = f_6(m,n)$.
\end{problem}

\begin{problem}
    Find $\frac{dy}{dx}$ if $y^3+x^3-xy=0$
\end{problem}

\section{Directional Derivatives}

\begin{problem}
Find $\mathcal{D}_{\hat{u}} f(1,1)$ where $f(x,y) = x^3-5y+3xy^2$ and $\hat{u} = \langle \frac{8}{17}, \frac{15}{17}\rangle $    
\end{problem}

\begin{problem}
Find the derivative of $f(x,y) = e^{x^2y}+\frac{y}{6x^3}$ in the direction that the angle $\frac{5\pi}{6}$ makes with the $x$-axis.
\end{problem}

\begin{problem}
Find the derivative of $f(x,y) = e^{x^2y}+\frac{y}{6x^3}$ in the direction of $\vec{w} = 2 \hat{i} - 9\hat{j}$.
\end{problem}

\section{Tangent Planes and Normal Lines}

\begin{problem}
Find the equation both the tangent plane and normal line to the following surfaces at the given points.

\begin{enumerate}
    \item $z = \frac{1}{2}x^2-\frac{1}{3}y^3+\frac{1}{4}x^4; (2,3, -3)$
    \item $z^2 + xyz = xy-2xz+yz+1 ; (1,2,1)$
    \item $f(x,y) = \sec{x}+\ln{\frac{xy}{\pi}}; (\frac{\pi}{3}, 3e,3)$
\end{enumerate}
\end{problem}


\section{Relative Extrema}

\begin{problem}
Find and classify all critical values of the following surface. \[f(x,y) = 7x-8y+2xy-x^2+y^3 \]
\end{problem}



\begin{problem}
Find and classify all critical values of the following surface. \[f(x,y)=(3x+4x^3)(y^2+2y)\]
\end{problem}

\begin{problem}
Find and classify all critical values of the following surface. \[f(x,y) = x^3 - 3x^2 + 3y^2 + y^3 - 5\]
\end{problem}


\begin{problem}
 Find the absolute maximum and minimum of $f(x,y) = 5x^2-3y^3+4y-x+1$ bounded by the region $R = \{(x,y) : 0\leq x\leq 4, 1\leq y\leq 8\}$   
\end{problem}
   
\begin{problem} 
    Find the absolute maximum/minimum and where they occur on the surface $f(x,y) = 2x^2+3xy+y^2$ bounded by the triangle with end points $(-4,-1),(3,-1),(1,2)$.
\end{problem}

\begin{problem}
Use Lagrange multipliers to maximize $f(x,y) = x^2 + 4y^2 - 2x +8y $ subject to the constraint: $x + 2y = 7$.
\end{problem}

\begin{problem}
Use Lagrange multipliers to maximize $f(x,y) = 8x^2 - 2y$ subject to the constraint: $x^2 + y^2 = 1$.
\end{problem}





\end{document} 